%% notes-tex/genesis/genesis-proposal-redacted/sections/objectives.tex
%% [[genesis-proposal-redacted]]
%% https://github.com/Mjvolk3/torchcell/tree/main/notes-tex/genesis/genesis-proposal-redacted/sections/objectives.tex
%% SOURCE: hand-transcribed from the proposal docx, "Project Objectives". The three
%% named industry partners are redacted.

\section{Project objectives\secstatus{todo}}\label{sec:objectives}

Building upon our previous accomplishments, we propose to develop a metabolic
engineering foundation models (MEFMs) through the following specific aims:

\begin{enumerate}
  \item Extend TorchCell to facilitate the development of MEFMs by expanding the
    TorchCell Database and improving metabolic engineering model architecture. We
    will extend the current TorchCell CGT to other organisms by finetuning and
    retraining on perturb-seq data\textsuperscript{12} collected from industrial
    platform organisms.
  \item Integrate an improved genome large language model (LLM) based on GenSLM for
    metabolic engineering (HiGenLM) finetuned on data from TorchCell Database.
  \item Use our biofoundries, including the iBioFAB, NSF iBioFoundry, SynBioAgent
    agentic AI system at the UIUC, the OPAL autonomous laboratory platform, and the
    IDeA co-scientist infrastructure at the ANL, to validate AI model predictions and
    generate large amounts of data in a high-throughput manner for finetuning and
    training these models.
\end{enumerate}

In Phase I (first 9 months), our main goal is to demonstrate proof of concept by
unifying the aforementioned AI models into an MEFM and establishing computational and
wet-lab experimental pipelines for scale-up studies in Phase II. In Phase II (next
three years), we aim to include all 100--250 industrially relevant and emerging
organisms used for metabolic engineering in the TorchCell Database and host metabolic
engineering benchmark datasets for the metabolic engineering community. We will also
collect data on at least twenty organisms widely used in the field of metabolic
engineering using our agentically driven biofoundries, and we will use TorchCell to
train metabolic engineering foundation models on at least 1000 GPUs in the
soon-to-be-established Equinox and Solstice systems at ANL. In addition, we will
collaborate with \redacted{two industry partners} to validate our MEFMs using their
platform organisms, and with \redacted{an instrument vendor} to enhance the technical
capabilities of our biofoundries for experimental validation and data generation (see
letters of support). Note that this project will not involve the construction,
alteration, maintenance and/or repair of public infrastructure in the United States.
